
\documentclass[../../../physics-standard_questions_all.tex]{subfiles}


\begin{document}

図のように，電池$\mathrm{E}$（起電力$V_0$），コンデンサ$\mathrm{C}_1$（容量$C$），$\mathrm{C}_2$（容量$2C$），
スイッチ$\mathrm{S}_1$，$\mathrm{S}_2$からなる回路がある．
はじめ，全てのスイッチは開いており，全てのコンデンサは帯電していない．

「$\mathrm{S}_1$を閉じて充分待ってから$\mathrm{S}_1$を開く．
続いて，$\mathrm{S}_2$を閉じて充分待ってから$\mathrm{S}_2$を開く」
という操作を$n$回くり返した後の$\mathrm{C}_2$の電荷を$y_n$とする
（$n=1,\:2,\:3,\cdots$）．

$y_\infty=\displaystyle\lim_{n \to \infty} y_n$を次の2通りの考え方で求めよ．

\begin{enumerate}

    \item 数列$\{y_n\}$の満たす漸化式を作り，それを解くことにより一般項を求める．

    \item 操作を充分に多数回くり返した後，もはや操作による電荷の移動がなくなることを前提とし，
        $y_\infty$が満たす関係式を作る．

\end{enumerate}

\vspace{0.2\baselineskip}
{\centering
    \includegraphics{\subfix{fig/fig_em_cx_04/fig_em_cx_04_00.pdf}}
\par}
\vspace{0.2\baselineskip}

\end{document}



